\section{Results}
\label{sec:results}

Results are structured around RQ1--RQ5.  For each question, the corresponding
estimand, family-level uncertainty, and policy or representation contrast are
presented together; held-out BBOB and external-suite estimates are kept
separate.

\subsection{RQ1: How often does the primary query have non-positive net utility?}
\label{sec:results-rq1}

RQ1 concerns the primary low-cost descriptor query under the main cost
weighting, $\lambda_T=1$ and $\lambda_M=0$. Query utility is decomposed into
the observed performance difference, downstream-selector regret, and measured
non-FE query cost. Evaluations consumed by query sampling are already absent
from the continuation budget of the query branch and are therefore not
subtracted again.

\begin{table*}[t]
\caption{State-level utility summary for the primary low-cost descriptor
query. Utility is larger-is-better, whereas final action loss is
smaller-is-better. Function families, rather than individual decision states,
form the independent units for uncertainty quantification.}
\label{tab:rq1-utility}
\centering
\scriptsize
\begin{tabularx}{\textwidth}{@{}l Y c c c c c Y@{}}
\toprule
Analysis stratum & Coverage & States & $\Pr(U_q\leq 0)$ & Mean $U_q$
& Median $U_q$ & Performance difference & Selector regret and query cost \\
\midrule
Overall & All prespecified BBOB-train families, dimensions, and decision
opportunities & -- & -- & -- & -- & -- & -- \\
By split & Separate BBOB-train family-OOF and frozen BBOB-validation estimates
& -- & -- & -- & -- & -- & -- \\
By dimension & 10D, 20D, and 40D & -- & -- & -- & -- & -- & -- \\
By decision stage & Each prespecified dynamic decision opportunity & -- & --
& -- & -- & -- & -- \\
\addlinespace
\multicolumn{8}{l}{\TBD{TBD-RQ1-TAB-01}: numerical estimates and
family-level intervals for all rows.} \\
\bottomrule
\end{tabularx}
\end{table*}

\begin{figure*}[t]
\centering
\fbox{\begin{minipage}{0.90\textwidth}
\centering
\TBD{TBD-RQ1-FIG-01}\\[0.8ex]
Distribution of $U_q$ relative to zero, with function-family-aware uncertainty
and separate views by dimension and actual FE ratio.
\end{minipage}}
\caption{Distributional analysis of utility for the primary query. The zero
reference distinguishes states in which the measured cost-adjusted utility is
non-positive.}
\label{fig:rq1-utility}
\end{figure*}

\paragraph{Statistical analysis.}
\TBD{TBD-RQ1-STAT-01} The family-balanced proportion of states with
$U_q\leq0$ is estimated with function-family-aware uncertainty; paired effects
decompose performance and cost contributions.  State counts are descriptive
and do not serve as independent inferential units.  Equivalence is evaluated
only under the prespecified clustering hierarchy, resample count, smallest
effect of interest, and primary interval or TOST rule.

\paragraph{Answer to RQ1.}
\TBD{TBD-RQ1-CLAIM-01} Whether the primary query has net benefit across the
evaluated decision states is undetermined.  The answer is conditional on the
estimated direction, effect size, interval, and observed state coverage, and is
restricted to the primary query rather than landscape analysis in general.

\FloatBarrier
\subsection{RQ2: Does algorithm-agnostic behavior predict query utility?}
\label{sec:results-rq2}

The active Decision Model candidates are LDA, Logistic Regression, and Ridge.
Their only model-family comparison uses the fixed B3 representation and nested
BBOB-train function-family out-of-fold (OOF) mean decision utility. AUROC,
Average Precision, and Spearman correlation are auxiliary measures, and
continuous-utility RMSE is defined only for Ridge. After this comparison, the
selected estimator family is fitted separately to B3 and T0, with a threshold
derived from the corresponding complete BBOB-train family-OOF scores.
BBOB-validation is evaluation-only and cannot influence preprocessing,
model-family choice, input-group comparison, or threshold estimation.

\begin{table*}[t]
\caption{Decision Model selection and evaluation structure for RQ2. Candidate
models are selected using B3 and nested BBOB-train family-OOF utility only.
The selected estimator family is subsequently fitted separately to B3 and T0;
BBOB-validation is reserved for evaluating these fixed choices without
refitting or reselection.}
\label{tab:rq2-prediction}
\centering
\scriptsize
\begin{tabularx}{\textwidth}{@{}l l Y c c c Y@{}}
\toprule
Analysis & Evaluation population & Estimator and representation
& Mean utility & AUROC / AP & Spearman / RMSE & Inferential role \\
\midrule
Model-family selection & BBOB-train nested family OOF & LDA with B3
& -- & -- & -- & Candidate comparison \\
Model-family selection & BBOB-train nested family OOF & Logistic Regression
with B3 & -- & -- & -- & Candidate comparison \\
Model-family selection & BBOB-train nested family OOF & Ridge with B3
& -- & -- & -- & Candidate comparison \\
\addlinespace
Representation comparison & BBOB-train family OOF & Selected estimator,
B3 versus T0 & -- & -- & -- & Paired B3--T0 effect \\
Frozen evaluation & BBOB-validation & Selected estimator with B3
& -- & -- & -- & Evaluation only \\
Frozen evaluation & BBOB-validation & Selected estimator with T0
& -- & -- & -- & Evaluation only \\
\addlinespace
\multicolumn{7}{l}{\TBD{TBD-RQ2-TAB-01}: model-selection, paired
representation, and frozen-evaluation estimates with family-level intervals.} \\
\multicolumn{7}{l}{RMSE is reported only when the estimator is Ridge; otherwise
it is not applicable. Frozen evaluation includes the paired B3--T0 effect.} \\
\bottomrule
\end{tabularx}
\end{table*}

\begin{figure*}[t]
\centering
\fbox{\begin{minipage}{0.90\textwidth}
\centering
\TBD{TBD-RQ2-FIG-01}\\[0.8ex]
(a) Outer-family OOF decision utility for LDA, Logistic Regression, and Ridge
with B3; (b) paired B3--T0 effects for the selected estimator family on
BBOB-train family-OOF predictions and on the separate frozen
BBOB-validation evaluation. Held-out families are the replication units.
\end{minipage}}
\caption{B3 model-family comparison and fixed B3--T0 evaluation. The
BBOB-validation panel has an evaluation-only role and cannot revise the model
choice made from BBOB-train utility.}
\label{fig:rq2-oof}
\end{figure*}

\paragraph{Statistical analysis.}
\TBD{TBD-RQ2-STAT-01} The paired outer-family BBOB-train OOF comparison for B3,
uses family-level effect sizes, confidence intervals, and Holm-corrected
candidate contrasts.  Conditional on that single model-family selection, the
same estimator family defines the B3--T0 comparison and the separate frozen
BBOB-validation evaluation.  Dependence-aware intervals use the prespecified
resampling unit and replication count; auxiliary score measures and
BBOB-validation cannot supersede the utility-based BBOB-train selection rule.

\paragraph{Answer to RQ2.}
\TBD{TBD-RQ2-CLAIM-01} Whether algorithm-agnostic behavior adds decision value
beyond FE ratio is undetermined and depends on the paired B3--T0 effects and
intervals.  BBOB validation is reserved for evaluating only the frozen models;
an interval that
does not support a stable difference yields an inconclusive answer.

\FloatBarrier
\subsection{RQ3: Does Decision-before-Feature reduce unhelpful query calls?}
\label{sec:results-rq3}

RQ3 retains all eight prespecified comparison roles but reports six distinct
outcome rows in the primary population.  The native path simultaneously
represents Never Query and the training-derived SBS, and the fixed
first-opportunity query path simultaneously represents Always Query and
Traditional AAS.  Random Analysis, the Time-only Controller, and
Decision-before-Feature provide the remaining online trajectories; VBS is a
per-problem hindsight reference.  Slash-labelled pairs are counted once in
summaries and inferential comparisons.  VBS is neither deployable nor the
statewise best observed action.

\begin{table*}[t]
\caption{Equal-budget RQ3 comparison.  Six distinct outcome rows represent all
eight prespecified comparison roles; slash-labelled rows combine definitions
that are analytically identical in the primary population and are counted once.
Panel A combines optimization performance with resource use. Panel B characterizes query-call quality,
including the frequency, within-call rate, and accumulated utility loss of
calls whose observed utility is non-positive.}
\label{tab:rq3-baselines}
\centering
\scriptsize
\begin{tabularx}{\textwidth}{@{}>{\raggedright\arraybackslash}p{0.21\textwidth} Y Y Y Y Y Y@{}}
\toprule
\multicolumn{7}{l}{\textit{Panel A: optimization performance and resource use}} \\
\addlinespace
Outcome row (comparison roles) & Mean decision utility & Final error & Runtime & Total FE & Query FE
& Query-call rate \\
\midrule
Never Query / SBS & 0 / N/A & -- & -- & -- & 0 / N/A & 0 / N/A \\
Always Query / Traditional AAS & -- & -- & -- & -- & -- & -- \\
Random Analysis & -- & -- & -- & -- & -- & -- \\
VBS & N/A & -- & N/A & -- & N/A & N/A \\
Time-only Controller & -- & -- & -- & -- & -- & -- \\
Decision-before-Feature & -- & -- & -- & -- & -- & -- \\
\bottomrule
\end{tabularx}

\vspace{0.8ex}
\begin{tabularx}{\textwidth}{@{}>{\raggedright\arraybackslash}p{0.21\textwidth} Y Y Y Y Y Y@{}}
\toprule
\multicolumn{7}{l}{\textit{Panel B: query-call quality}} \\
\addlinespace
Policy & $U_q\leq0$ calls & Within-call rate & Accumulated utility loss
& Beneficial-state capture & Utility capture & Precision \\
\midrule
Always Query / Traditional AAS & -- & -- & -- & -- & -- & -- \\
Random Analysis & -- & -- & -- & -- & -- & -- \\
Time-only Controller & -- & -- & -- & -- & -- & -- \\
Decision-before-Feature & -- & -- & -- & -- & -- & -- \\
\addlinespace
\multicolumn{7}{@{}>{\raggedright\arraybackslash}p{\textwidth}@{}}{%
\TBD{TBD-RQ3-TAB-01}: numerical estimates and matched family-level intervals
for Panels A and B.} \\
\multicolumn{7}{@{}>{\raggedright\arraybackslash}p{\textwidth}@{}}{%
Never Query / SBS and VBS are omitted from Panel B because call-conditional
quantities are undefined for static references and Never Query makes no calls.} \\
\multicolumn{7}{@{}>{\raggedright\arraybackslash}p{\textwidth}@{}}{%
A slash-labelled row denotes one shared outcome, not an average of two
estimates; it enters each applicable analysis once.} \\
\bottomrule
\end{tabularx}
\end{table*}

\begin{figure*}[t]
\centering
\fbox{\begin{minipage}{0.90\textwidth}
\centering
\TBD{TBD-RQ3-FIG-01}\\[0.8ex]
Final optimization loss, measured runtime, and query-call rate for each
distinct primary-population outcome under equal FE budgets, with function-family-level
uncertainty rather than a frontier based only on point estimates.
\end{minipage}}
\caption{Cost--performance comparison among the equal-budget outcomes.
Coincident comparison roles are shown once.}
\label{fig:rq3-cost-performance}
\end{figure*}

\paragraph{Statistical analysis.}
\TBD{TBD-RQ3-STAT-01} The matched function-family Friedman comparison,
Holm-adjusted paired Wilcoxon contrasts, effect sizes, and confidence intervals
jointly cover mean decision utility as defined in
Eq.~\eqref{eq:policy-decision-utility},
final error, runtime, and both the rate and accumulated cost of calls with
$U_q\leq0$.  The dependence-aware interval construction and equivalence rule
are specified independently of the observed policy differences.
Omnibus and follow-up tests use each distinct applicable outcome once; N/A
roles are omitted endpoint by endpoint, and neither coincident pair is
duplicated.

\paragraph{Answer to RQ3.}
\TBD{TBD-RQ3-CLAIM-01} Whether Decision-before-Feature reduces calls with
non-positive observed utility is undetermined.  A supported answer requires an
interval-supported reduction together with the equal-budget effects on final
optimization performance and runtime; call rate alone is insufficient.

\FloatBarrier
\subsection{RQ4: Does the frozen controller generalize beyond BBOB training families?}
\label{sec:results-rq4}

RQ4 reserves the fully frozen procedure for held-out BBOB-validation families,
CEC2017, CEC2022, and an engineering-problem evaluation to be specified
prospectively. None
of these evaluation sets can refit the Selection Reference, Decision Model,
preprocessing, representation, or OOF utility threshold. Coverage and query
extraction failures are part of the suite-specific evaluation rather than
grounds for silently excluding problems.

\begin{table*}[t]
\caption{Suite-specific evaluation of the frozen controller for RQ4. Coverage
and query-extraction failure rates accompany utility and final optimization
performance so that estimates remain conditional on observed evaluation
coverage. Paired effects compare Decision-before-Feature with Never Query,
Always Query, and the Time-only Controller, in that order.}
\label{tab:rq4-generalization}
\centering
\scriptsize
\begin{tabularx}{\textwidth}{@{}l Y c c c c Y@{}}
\toprule
Evaluation set & Scope & Problem--dimension--seed coverage & Mean utility & Final error
& Query failures & Paired effect and interval \\
\midrule
BBOB-validation & Held-out BBOB function families & -- & -- & -- & -- & -- \\
CEC2017 & Cross-benchmark evaluation & -- & -- & -- & -- & -- \\
CEC2022 & Cross-benchmark evaluation & -- & -- & -- & -- & -- \\
Engineering problems & Problem set to be specified prospectively & -- & -- & -- & -- & -- \\
\addlinespace
\multicolumn{7}{l}{\TBD{TBD-RQ4-TAB-01}: suite-specific coverage,
performance, failure, and matched-effect estimates.} \\
\bottomrule
\end{tabularx}
\end{table*}

\begin{figure*}[t]
\centering
\fbox{\begin{minipage}{0.90\textwidth}
\centering
\TBD{TBD-RQ4-FIG-01}\\[0.8ex]
Suite-specific paired effects and confidence intervals for the frozen
Decision-before-Feature controller relative to Never Query, Always Query, and
the Time-only Controller. Held-out BBOB families are shown separately from
CEC and engineering evaluations.
\end{minipage}}
\caption{Held-out-family and cross-benchmark evaluation of the frozen
controller. Separate panels preserve differences in benchmark scope and data
coverage.}
\label{fig:rq4-generalization}
\end{figure*}

\paragraph{Statistical analysis.}
\TBD{TBD-RQ4-STAT-01} Suite-specific problem- or function-family-level paired
effects, confidence intervals, and multiplicity-adjusted benchmark contrasts
evaluate the frozen policy.  Query-extraction failure sensitivity is separate
from the primary comparison, and heterogeneous effect directions or coverage
preclude a pooled generalization claim.

\paragraph{Answer to RQ4.}
\TBD{TBD-RQ4-CLAIM-01} Held-out BBOB performance and cross-benchmark
generalization are undetermined.  Any eventual answer is suite-specific,
distinguishes BBOB validation from CEC2017, CEC2022, and engineering
evaluations, and remains conditional on query-extraction coverage.

\FloatBarrier
\subsection{RQ5: Which behavior groups are associated with query utility, and when?}
\label{sec:results-rq5}

RQ5 reuses the estimator family selected with B3 in RQ2. That estimator family
is fitted separately to T0, B1, B2, and B3 using identical samples,
function-family splits, train-only preprocessing, and an input-specific
BBOB-train family-OOF threshold. The ablation neither repeats model-family
selection nor changes the primary B3 representation. T0 contains only FE
ratio; B1, B2, and B3 contain the prespecified 19-, 25-, and 31-input
permutation-invariant behavior representations. Coefficients, discriminant
directions, and associations characterize the fitted rule and are not causal
effects.

\begin{table*}[t]
\caption{Prespecified input-group ablation for RQ5. The estimator family is
the one selected with B3 in RQ2 and is fitted separately to each representation.
All groups retain their own BBOB-train family-OOF threshold; the comparison
does not select a replacement for B3.}
\label{tab:rq5-ablation}
\centering
\scriptsize
\begin{tabularx}{\textwidth}{@{}l Y c c c c Y@{}}
\toprule
Group & Added information & Inputs & Mean utility & AUROC / AP
& Spearman / Ridge RMSE & Paired difference from T0 \\
\midrule
T0 & Optimization stage (FE ratio) only & 1 & -- & -- & -- & -- \\
B1 & Progress, diversity, set and fitness-distribution change, stagnation,
convergence, and elite concentration & 19 & -- & -- & -- & -- \\
B2 & B1 plus longitudinal centroid, covariance, fitness-spread, and diversity
recovery summaries & 25 & -- & -- & -- & -- \\
B3 & B2 plus set-motion proxies and the three operational maturity variables
& 31 & -- & -- & -- & -- \\
\addlinespace
\multicolumn{7}{l}{\TBD{TBD-RQ5-TAB-01}: paired family-level ablation
estimates, auxiliary scores, and intervals for all four representations.} \\
\multicolumn{7}{l}{The paired difference from T0 concerns decision utility.
RMSE is not applicable unless the RQ2-selected estimator is Ridge.} \\
\bottomrule
\end{tabularx}
\end{table*}

\begin{figure*}[t]
\centering
\fbox{\begin{minipage}{0.90\textwidth}
\centering
\TBD{TBD-RQ5-FIG-01}\\[0.8ex]
Fold-stable standardized coefficients for Logistic Regression or Ridge, or
standardized discriminant directions for LDA, together with the prespecified
association between behavior-derived search maturity and observed query
utility. Uncertainty and direction changes are separated by family, dimension,
and query configuration.
\end{minipage}}
\caption{Behavior interpretation and search-maturity--utility analysis.
Displayed coefficients or discriminant directions describe the fitted
decision score and do not imply causal effects.}
\label{fig:rq5-interpretation}
\end{figure*}

\paragraph{Statistical analysis.}
\TBD{TBD-RQ5-STAT-01} The paired outer-family effects and
multiplicity-adjusted contrasts compare T0/B1/B2/B3 for the single estimator
family selected in RQ2.  Interpretation also depends on the sign and magnitude
stability of standardized coefficients
or discriminant directions across outer folds and families, Spearman
associations with observed utility, uncertainty around the fitted
maturity--utility curve, and the prespecified dimension and query-configuration
analyses.  Visual shape alone is insufficient evidence of non-monotonicity,
and the ablation cannot repeat model or representation selection.

\paragraph{Answer to RQ5.}
\TBD{TBD-RQ5-CLAIM-01} Whether any behavior increment is informative beyond
optimization stage is undetermined.  Such an interpretation requires agreement
among the paired ablation, parameter or discriminant stability, and held-out
analyses.  Representation- or family-dependent effects instead imply a
correspondingly restricted answer; none can replace B3 through post hoc
selection.

\FloatBarrier
\subsection{Prespecified sensitivity analyses and diagnostics}
\label{sec:results-sensitivity}

These analyses qualify the RQ1--RQ5 conclusions without defining an additional
research question or changing the primary target. Each alternative query
configuration traverses the same complete query-specific sequence of action
selection, utility estimation, decision modelling, and frozen evaluation as
the primary query. The non-primary $\lambda_T$ values do not replace
$\lambda_T=1$. Action-relation indicators define descriptive strata only, and
downstream-selector variants neither enter the Decision Model nor replace the
policy baselines.

\begin{table*}[t]
\caption{Prespecified sensitivity analyses and downstream-selector
diagnostics. Interpretation is restricted to the listed levels and cannot
revise the primary cost weighting, input representation, or policy comparison.}
\label{tab:sensitivity-diagnostics}
\centering
\scriptsize
\begin{tabularx}{\textwidth}{@{}l Y Y Y@{}}
\toprule
Analysis & Prespecified levels & Quantities & Interpretation boundary \\
\midrule
Query configuration & Primary low-cost descriptor query, standard
\textit{pflacco} query, and broad \textit{pflacco} query & Query-specific
action selection, utility, decision-model, policy, and frozen-evaluation
estimates & Agreement indicates robustness only across these three
configurations; disagreement indicates representation or cost dependence \\
$\lambda_T$ sensitivity & $0$, $0.25$, $0.5$, $1$, and $2$, with
$\lambda_M=0$ & Utility distributions, policy summaries, effect sizes, and
intervals & $\lambda_T=1$ remains the primary conclusion basis \\
Action-relation strata & Selected action equals the default action; selected
action equals the prefix algorithm; population handoff is required & Counts,
utility components, selector regret, and policy effects & Descriptive
stratification only; no action-relation indicator enters the Decision Model \\
Downstream-selector inputs & State information only, query descriptors only,
and their combination; remaining budget is included in every variant &
Cross-fitted action-loss regression performance, selected-action loss, and
selector regret & Downstream-selector diagnosis only; these variants are not
Decision Model candidates or substitutes for policy baselines \\
\addlinespace
\multicolumn{4}{l}{\TBD{TBD-SENS-DIAG-01}: estimates and intervals for all
prespecified sensitivity and diagnostic comparisons.} \\
\bottomrule
\end{tabularx}
\end{table*}

\FloatBarrier
